(1) is Theorem \claimref{standard-model.grip-mass-two-rung-module}(ii),(iii), read on the species of the generation under (S) (\claimref{standard-model.hyp-s-slot-statistics}): the mass-squared is a quadratic form on each sector, identically zero on the slot-singlet sectors and the area on the slot-doublet sectors; it is invariant under the rung mixing by Theorem \claimref{standard-model.moment-map-grip-mass-are-insensitive} but not under the pair mixing (Proposition \claimref{standard-model.frozen-point-pair-mixing}), which is why the doublet entries are forms on their sectors rather than scalars, while the zeros are exact whatever the pair state. (2) is Corollary \claimref{standard-model.under-s-frozen-points-split-inverse}. (3) is Proposition \claimref{standard-model.under-s-bracket-direction-charged-not} with Proposition \claimref{standard-model.mass-bracket-charged-singlet}'s caveat that the bracket of an excited composite is not defined. (4): the objects the series defines that carry mass are the grip and the bracket, Fact \claimref{standard-model.imported-grip-area-casimir-form} and Proposition \claimref{standard-model.mass-bracket-charged-singlet}; (1) to (3) exhaust what they do on the module; the six numbers are otherwise free by Proposition \claimref{standard-model.invariant-operators-module}. The identification of the missing sector with the Higgs doublet's origin is Observation \claimref{standard-model.configuration-sector-carries-higgs-doublets-quantum}'s first clause stated as a problem, not as a result: what is proved is that the split is not derived in the series.
